\documentclass[11pt]{article}
\usepackage{fontspec,amsmath,amssymb,amsthm,hyperref}
\usepackage[margin=1in]{geometry}
\setlength{\emergencystretch}{3em}
\title{Uniform potential gaps and marginal-center displacement}
\author{LangevinLean mathematical exposition}
\date{}
\begin{document}\maketitle
\noindent Source module: \texttt{ProximalLangevin/Perturbation.lean}. Compile this proof document with XeLaTeX or LuaLaTeX. The accompanying HTML places the complete Lean declarations beside every proof.


Let \(V\le U\le V+\delta\), with \(V\) differentiable and \(m\)-strongly convex, \(m>0\), and \(U\) continuous. Write \(k_W(\theta)=\int e^{-W(\theta,x)}dx\) and \(F_W=-\log k_W\). This module derives positive finite marginals and actual maximizers before comparing centers.


\section{Quadratic growth needs only strong convexity}


For a globally minimizing \(z\) of an \(m\)-strongly convex function \(F\), apply strong convexity to \((1-t)z+tx\) and use \(F(z)\le F((1-t)z+tx)\). For \(t>0\), divide the resulting inequality by \(t\), then let \(t\downarrow0\). This gives \(F(x)-F(z)\ge(m/2)\|x-z\|^2\) without differentiating \(F\).



\noindent Lean declarations: \texttt{strongConvex\_minimum\_growth}.

\section{Integrate the gap before taking logarithms}


For any potentials \(V\le U\le V+\delta\), exponentiation and integration in the latent variable give


\[
e^{-\delta}k_V(\theta)\le k_U(\theta)\le k_V(\theta),
\qquad k_W(\theta)=\int e^{-W(\theta,x)}dx.
\]


The smooth strongly convex comparison potential \(V\) supplies positive finite integrals and an integrable Gaussian bound. The original marginal is continuous and decays at infinity; hence it has an actual maximizer. Taking negative logarithms, which is justified by positivity, gives


\[
F_V(\theta)\le F_U(\theta)\le F_V(\theta)+\delta,
\qquad F_W=-\log k_W.
\]


Let \(\theta_U\) minimize \(F_U\) and \(\theta_V\) minimize the \(m\)-strongly convex \(F_V\). Minimality and the sandwich imply \(F_V(\theta_U)-F_V(\theta_V)\le\delta\). Strong-convexity growth about its minimizer therefore yields


\[
\frac m2\|\theta_U-\theta_V\|^2\le\delta.
\]


This argument never differentiates the original marginal and never integrates an asymptotic remainder.



The sharp quadratic-growth lemma uses strong convexity and global minimality, without differentiability of the marginal potential. The comparison model supplies an integrable Gaussian bound for the original weights, giving continuity and existence of an original marginal maximizer by dominated convergence and decay at infinity. Finite positive marginal integrals are established before logarithms and real conversions.



\noindent Lean declarations: \texttt{marginalDensity\_sandwich}, \texttt{marginalPotential\_sandwich}.

\section{Establish original-marginal existence}


The inequality \(e^{-U}\le e^{-V}\) and strong-convexity Gaussian bounds prove integrability. Continuity follows by dominated convergence on parameter neighborhoods, using the comparison potential to provide an integrable latent majorant. The same Gaussian estimate proves decay at parameter infinity. A positive continuous function that vanishes at infinity attains its maximum on a sufficiently large compact ball.



\noindent Lean declarations: \texttt{integrable\_original\_marginal}, \texttt{continuous\_original\_marginal}, \texttt{exists\_original\_marginalMaximizer}.

\section{Strong growth controls the center shift}


Positive finite marginal densities make \(-\log\) order-reversing, so an original marginal maximizer minimizes \(F_U\). At minimizers \(\theta_U,\theta_V\), the sandwich and minimality give \(F_V(\theta_U)\le F_U(\theta_U)\le F_U(\theta_V)\le F_V(\theta_V)+\delta\). The preserved marginal strong convexity yields \(\|\theta_U-\theta_V\|^2\le2\delta/m\); taking the nonnegative square root gives the distance estimate.



\noindent Lean declarations: \texttt{original\_maximizer\_minimizes}, \texttt{marginal\_center\_distance\_sq}, \texttt{marginal\_center\_distance}.

\end{document}
